\null\vspace{0.1cm}
%~ \begin{eng}
%~ \centerline{\Chapter{Introduction}}
%~ \end{eng}
%~ \vskip 7pt

\begin{nohyp}
\begin{eng}
\centerline{\Section{Royal House of Mysore Yoga Manuscript}\label{introsec1}}
\end{eng}
\vskip 22pt

\begin{multicols*}{2}
%\begin{nohyphens}
\begin{eng}
\lettrine[lines=3,lhang=0.04,nindent=0em]{T}{his volume}, the second publication in the \emph{Chitsthithananda} series, brings to light a literary treasure preserved in one of India's most influential royal courts, the Royal House of Mysore. In the private collection of the Royal House of Mysore in the custody of the head of the family, H.\,H.\ Maharani Dr.\ Pramoda Devi Wadiyar, is a manuscript that stands among the most significant surviving documents for the study of India's yoga traditions. Designated the \emph{Royal House of Mysore Yoga Manuscript},\endnote{Throughout this volume, for the sake of brevity, the codex will be referred to as the \emph{Mysore Yoga Manuscript}.} this exquisitely illustrated artefact dates to the early nineteenth century and was likely produced during the reign of His Highness Krishnaraja Wadiyar III (also known as Śrī Mummaḍi Kṛṣṇarāja Wadiyar), the twenty-second Maharaja of Mysore~(r.~1794--1868~CE).

Deeply rooted in India's tradition of yoga yet flavoured with regionally specific nuances, the yoga texts preserved in this manuscript figure large in the history of yoga in Mysore and in the Royal family's patronage of the practice of yoga postures (\emph{yogāsana}). The manuscript is also an essential source for understanding the broader history of physical yoga in India and its subsequent global influence.

This volume presents a scholarly facsimile edition of three yoga manuals (\emph{paddhati}) preserved in the \emph{Mysore Yoga Manuscript.} It comprises high-resolution colour facsimiles, a critically edited transcription of the Sanskrit text in Devanagari script, an annotated English translation, and a Kannada translation. The three manuals are attributed to the perfected yogi~(\emph{siddha}) named Koraṇṭaka, who is said to have been a disciple of Gorakṣanātha. Gorakṣanātha is traditionally regarded as one of the founders of Haṭhayoga (``the yoga of force''), a type of yoga attained through physical means that emerged at the beginning of the second millennium. Koraṇṭaka's manuals, which we designate \emph{Koraṇṭakapaddhati}~\emph{I}, \emph{II}, and \emph{III}, respectively, were likely composed in eighteenth-century Maharashtra and later brought to Mysore, where they were preserved in Kannada script and beautifully illustrated in the early nineteenth century.

This volume also includes a comprehensive introduction: Section~\ref{introsec1} discusses the date, physical features, script, and contents of the \emph{Mysore Yoga Manuscript} and situates the manuscript within the corpus of works produced during the reign of Krishnaraja Wadiyar III. Section~\ref{introsec2} examines what is known about Koraṇṭaka, the yoga tradition associated with him, and the manuals attributed to him, with particular attention to their provenance and historical significance. Section~\ref{introsec3} turns to the sublime artwork of \emph{Koraṇṭakapaddhati}~\emph{I}, considering its artistic significance and contribution to the Mysore school of painting. Section~\ref{introsec4}, demonstrates how sustained royal patronage of the teaching and practice of \emph{yogāsana} at the Mysore royal court laid the foundations for the subsequent development of yoga within India and beyond, thereby contributing to the emergence of modern postural yoga. The final section, Section~\ref{introsec5}, sets out the editorial conventions underlying the facsimile and Devanagari editions, as well as the English and Kannada translations.

The \emph{Mysore Yoga Manuscript} contains several other important works on Haṭhayoga that fall beyond the scope of this publication, including a unique illustrated version of the fifteenth-century \emph{Haṭhapradīpikā}, which has been augmented with additional material. It also preserves a short Sanskrit work on \emph{vajrolīmudrā} entitled the \emph{Kāminīkalpa}, as well as an illustrated Hindi-inflected vernacular work entitled the \emph{Jñānaratnamālā}. The latter two works are unknown to modern scholarship on yoga. It is hoped that facsimile editions of these works will be published in future volumes of the same series, further attesting to the rich religious, intellectual, and artistic legacy of the Mysore kingdom.

The historical significance of the \emph{Mysore Yoga Manuscript} is manifold. It is one of only a few known illustrated works of yoga and a rare example of early Mysore-style painting executed by court artists during the reign of Krishnaraja Wadiyar III. Above all, however, it provides tangible evidence of the Royal House of Mysore's sustained patronage of yoga. Annotations in the manuscript reveal that it served as the source for an entry on the number 80 in the \emph{Saṅkhyāratnamālā} (``The Jewel Garland of Numbers'') and the section on yoga in the \emph{Śrītattvanidhi} (``The Illustrious Treasures of Reality''), two celebrated compendiums commissioned by Krishnaraja Wadiyar III. Richly illustrated, the \emph{Śrītattvanidhi} celebrates the arts, culture, and religion of the Mysore kingdom and, through its chapter on yoga, demonstrates the importance accorded to \emph{yogāsana}s since the early nineteenth century.

As we discuss in Section~\ref{introsec4.3} of the Introduction, the \emph{Koraṇṭakapaddhati} subsequently acquired a mythical status as the lost \emph{Yogakuraṇṭa} (or \emph{Yogakuraṇṭi}) in the imagination of practitioners within the lineage of Tirumalai Krishnamacharya (18 November 1892--28~February 1989). T.~Krishnamacharya, who served as the principal yoga teacher at the Royal House of Mysore from 1931 to \emph{circa} 1950 (Sjoman~1996:~50), privately kept a book of illustrations reproduced from the \emph{Haṭhapradīpikā} and \emph{Koraṇṭakapaddhati~I} of the \emph{Mysore Yoga Manuscript.} He was undoubtedly influenced by its sequences of \emph{yogāsana}s, as evinced by his citation of the \emph{Yogakuraṇṭi} in his book \emph{Yogāsanagaḷu} (``Yoga Postures''), published in~1941.

Through the patronage of the Royal House of Mysore, Krishnamacharya's system of \emph{yogāsana}s has had a profound and lasting influence on contemporary yoga, chiefly by way of his students who disseminated his teachings to an international audience. Today, practitioners in cities around the world, from London, Paris, New York, and Shanghai to Tokyo, Sydney, Moscow, Cape Town, and Delhi, can be seen practising dynamic forms of postural yoga that owe much to the tradition represented by Koraṇṭaka's manuals in the \emph{Mysore Yoga Manuscript}.

It is our great honour to present the teachings of the illustrious sage Koraṇṭaka. We hope this volume will inspire readers, practitioners, teachers, and researchers of yoga and art history around the world, and that they will delight in this remarkable treasure of the Royal House of Mysore.
\vskip 18pt

\Subsection{Date}\label{introsec1.1}
\vskip 9pt

\noindent
The \emph{Mysore Yoga Manuscript} was likely produced between 1815 and 1849~CE, with the balance of the evidence favouring a date closer to the beginning of this interval. The upper limit of the date (1849~CE) is established by the completion of the \emph{Saṅkhyāratnamālā}, which drew upon the \emph{Haṭhapradīpikā} and \emph{Koraṇṭakapaddhati~I} of the \emph{Mysore Yoga Manuscript} (Section~\ref{introsec1.5}). The lower limit (1815~CE) corresponds with the beginning of the literary programme commissioned by Krishnaraja Wadiyar III.

It is generally believed that Krishnaraja Wadiyar III's patronage of literature and the arts began in earnest after the `Nagara Peasant Rebellion' (1830--1831~CE), when his administrative authority and control over land revenue were transferred to the British (Ikegame 2013: 56--57, Simmons 2020: 135). Nevertheless, several lines of evidence suggest that his literary patronage began earlier. For example, a lithic inscription dated 1829~CE in the Prasannakṛṣṇaswāmī Temple in the Mysore Fort records Krishnaraja Wadiyar III's service in the form of nine jewels (\emph{navaratna}). The ninth, the jewel of learning (\emph{sārasvataratna}), commemorates his patronage of poetical works and Kannada-language commentaries on texts such as the \emph{Bhāgavatapurāṇa}, \emph{Adhyātmarāmāyaṇa}, and other puranic and epic literature.\endnote{See Rao 1930: 2962, Simmons 2019: 209. The inscription is transliterated and translated in \emph{Institute of Kannada Studies} 1976: 181--183, 765--766. We wish to thank Caleb Simmons for sending us the reference in the \emph{Institute of Kannada Studies} volume.} Moreover, his patronage had already produced important literary works prior to this time, including the \emph{Devīmahātymasaptaśatī} (1815~CE), \emph{Śaṅkarasaṃhitā} (1820~CE) and \emph{Lalitopākhyāna} (1823~CE). The first volumes of the \emph{Śrītattvanidhi} may also have been produced as early as 1823~CE (Hāvanūra 2011: 178).

\begin{figure*}[tb]
\centering
\includegraphics[width=\textwidth]{intro/01.jpg}
\caption{Label on the Spine of the \emph{Mysore Yoga Manuscript}.}\label{introfig1}
\end{figure*}

Because the \emph{Mysore Yoga Manuscript} served as a source for both the \emph{Saṅkhyāratnamālā} and the \emph{Śaivanidhi} of the \emph{Śrītattvanidhi}, it must predate both works (see Section~\ref{introsec1.5}). This provides strong evidence for placing its production in the early phase of Krishnaraja Wadiyar III's literary programme, when manuscripts were scribed in Kannada script and archived as source materials for his large compendiums.



Additional evidence that the \emph{Mysore Yoga Manuscript} was produced for Krishnaraja Wadiyar~III's personal library is provided by a label attached to the spine of its hand-stitched cloth binding (Figure~\ref{introfig1}). The remaining legible portion of the label appears to identify its contents as \emph{yogāsana}s (\emph{yogāsanagaḷu}).\endnote{The discernible letters on the label are: \emph{yogāsa}[\emph{na}]\emph{gaḷu} 80 \textbar\textbar{} \emph{hetvā}++\emph{ḷu} \emph{ālavū}+++. We wish to thank Viswanatha Gupta for his help with this reading.} More significantly, the label is decorated with a floral border framed by red lines and an inner teal-coloured line. Identical decorative features appear on the label attached to the spine of the illustrated \emph{Bhāgavatapurāṇa} (\emph{circa} 1820--1840~CE) produced by Krishnaraja Wadiyar III's court. Embossed with the Maharaja's royal seal and bound in a red silk cover (Goswami et al. 2019: Figs. 3, 4, and~6), this illustrated \emph{Bhāgavatapurāṇa} contains artwork similar in style to that of the now dispersed \emph{Devīmahātyma} (\emph{circa} 1815~CE), which was also produced early in Krishnaraja Wadiyar III's literary programme. The close similarity between the two labels suggests that the \emph{Mysore Yoga Manuscript} was housed in the same royal library.
\columnbreak



\Subsection{Physical Features, Script, and Layout}\label{introsec1.2}
\vskip 9pt

\noindent
The \emph{Mysore Yoga Manuscript} was produced in the conventional South Asian, horizontal \emph{pothi} format, with both the text and illustrations arranged in landscape orientation. The handmade-paper folios were cut in half and bound along the upper long edge, forming a long, narrow rectangular manuscript measuring 32\,cm in width and 14\,cm in height. Each text in the manuscript appears to have been originally copied onto folios of slightly different sizes. When these folios were gathered and bound into a single bundle, they were trimmed to create uniform edges, resulting in the partial loss of some marginal notes. These features indicate that the \emph{Mysore Yoga Manuscript} is a composite codex comprising more than a dozen manuscripts, copied separately in black and red ink and later assembled, perhaps because of their related subject matter on yoga. Its composite nature accounts for the variations in the cream- and green-coloured handmade papers and in the handwriting throughout the bundle.

\begin{center}
\includegraphics[width=\columnwidth]{intro/02.png}
{Figure 2: Selected letters from the \emph{Mysore Yoga Manuscript}.}
\refstepcounter{figure}\label{introfig2}
\end{center}

Overall, the script accords with the descriptions of Kannada presented by Kittel (1908) and Grünendahl (2001), although certain forms may reflect the idiosyncrasies of the scribe's hand. However, the following features occur with notable regularity. As shown in Figure~\ref{introfig2}, the letters \emph{gha}, \emph{ca}, \emph{ta}, and \emph{ṣa} often more closely resemble their Telugu counterparts. Likewise, \emph{sa} more closely resembles its Telugu counterpart, although unlike its standard forms in either Kannada or Telugu, the headmark and body of \emph{sa} are often joined and written with a single stroke. Similarly, although \emph{pa} is typically formed with two strokes, its headmark and body are sometimes joined. As in the Telugu script, \emph{gha}, \emph{pa}, \emph{ṣa}, and \emph{sa} lack the dot characteristic of their standard Kannada forms.




Other recurring scribal features include the substitution of the letter \emph{ḷa} for \emph{la}; the use of a variant form of \emph{bha} resembling \emph{kha} but distinguished by an extended upper-right stroke; and the inconsistent use of the vertical stroke that distinguishes the aspirated letters \emph{dha} and \emph{bha}. Finally, when \emph{va} forms part of a conjunct consonant (e.g., \emph{tva}, etc.), its conjunct form differs from those found in both the Kannada and Telugu scripts.

The illustrated folios of \emph{Koraṇṭakapaddhati~I} share a consistent layout in several respects. As seen in Figure~\ref{introfig4}, each folio presents four paintings, one in each quadrant, with the textual description of the \emph{āsana} placed above the rectangular frame enclosing the image. Each frame is formed by double red-ink lines, with a secondary horizontal double line beneath the figure. However, the size of each frame varies depending on whether the yogi's body is oriented horizontally or vertically, with frames measuring, on average, 12 cm in width and 9 cm in height. Particular attention was given to the compositional arrangement of each illustration and its relationship to the others, as discussed in Section~\ref{introsec3}.
\vskip 18pt

\Subsection{Content}\label{introsec1.3}
\vskip 9pt

\noindent
The \emph{Mysore Yoga Manuscript} consists of ninety-two paper folios, and its front cover is now lost. This composite codex was most likely scribed, illustrated, and assembled in Mysore (see Section~\ref{introsec1.1}). In view of its place of compilation and the fact that five of its principal texts concern yoga and together comprise nearly ninety per cent of its folios, it is referred to here as the \emph{Mysore Yoga Manuscript} (designated M in endnotes of the edition).

\begin{table*}[tb]
\centering
\renewcommand{\arraystretch}{1.6}
\caption{Contents of the \emph{Mysuru Yoga Manuscript}}\label{introtable1}
\vskip 6pt

\begin{tabular}{lcc}
\hline
\rowcolor{tableshade}
Name of the Text & Total No. of Folios & Cumulative Folio Numbers \\
\hline
\emph{Haṭhapradīpikā} & 23.5 & 1--23.5 \\
\emph{Kāminīkalpa} & 1 & 23.5--24.5 \\
\emph{Haṭhapradīpikā} (cont'd) & 8 & 24.5--32.5 \\
\rowcolor{tableshade}
\emph{Koraṇṭakapaddhati I} & 39.5 & 32.5--72 \\
{[}Fragment 1{]} & 0.5 & 72--72.5 \\
\rowcolor{tableshade}
\emph{Koraṇṭakapaddhati II} & 3.5 & 72.5--76 \\
\rowcolor{tableshade}
\emph{Koraṇṭakapaddhati III} & 1 & 76--77 \\
{[}Fragment 2{]} & 1.5 & 77--78.5 \\
\emph{Deśikastotra} & 2.5 & 78.5--81 \\
{[}Fragment 3{]} & 0.5 & 81--81.5 \\
\emph{Śivanṛtya} & 3 & 81.5--84.5 \\
\emph{Yantradvāra} (in the \emph{Ākāśakalpa}) & 0.5 & 84.5--85 \\
{[}Fragment 4{]} & 1 & 85--86 \\
\emph{Jñānaratnamālā} & 6 & 86--92 \\
\hline
\rowcolor{tableshade}
Total No. of Folios & 92 & 92 \\
\hline
\end{tabular}
\end{table*}

As shown in Table~\ref{introtable1}, it contains nine individual texts and four untitled text fragments that lack colophons. The first work is a finely illustrated version of the \emph{Haṭhapradīpikā}, comprising five chapters and containing substantial additional material. The first, third, fourth, and fifth chapters begin on new folios with the invocation ``Homage to the glorious Gaṇeśa'' (\emph{śrīgaṇeśāya namaḥ}), as though each were conceived as an independent text. Although the work is arranged in five chapters, the colophons of the final two designate each chapter as the fourth. Both concern Rājayoga and appear to correspond to the first and second halves of the fourth chapter in the early recensions of the \emph{Haṭhapradīpikā}.

The first chapter of the \emph{Haṭhapradīpikā} is embellished with twenty-five watercolour paintings of yogis performing a selection of \emph{āsana}s. Although the section on \emph{āsana} opens with a list of the fifteen postures found in most recensions of the \emph{Haṭhapradīpikā,} this unique recension incorporates additional material, including descriptions of the following yogic postures: \emph{mṛgasvastika}, \emph{cakrāsana}, \emph{prasāritāsana}, \emph{daṇḍāsana}, \emph{kapālāsana}, \emph{vajrāsana}, \emph{yogāsana}, \emph{añjalikāsana}, \emph{ardhacandrāsana}, \emph{yogapaṭṭāsana}, and \emph{garuḍāsana.} The descriptions of several of these additional \emph{āsana}s, including \emph{cakrāsana} and \emph{kapālāsana,} correspond to those in the long recension of the \emph{Yogacintāmaṇi} (Birch 2024). The author of that recension attributes these \emph{āsana}s to the early Śaiva work known as the \emph{Dharmaputrikā} (Barois 2020).

\begin{figure*}[tb]
\centering
\includegraphics[width=\textwidth]{intro/03.jpg}
\caption{Fragment 1 of the \emph{Mysore Yoga Manuscript}, folio \emph{72 verso}.}\label{introfig3}
\end{figure*}

The third chapter of the \emph{Haṭhapradīpikā} includes an additional folio of twenty-five verses immediately following the section on \emph{vajrolīmudrā}. The insertion begins with the invocation, ``Homage to the glorious Gaṇeśa'' (\emph{śrīgaṇeśāya namaḥ}), and concludes with the statement, ``Thus, the \emph{Kāminīkalpa} is complete'' (\emph{iti kāminīkalpaḥ samāptaḥ}). It appears to constitute a short work on \emph{vajrolī} intended for householders and contains several features characteristic of Koraṇṭaka's teaching on this \emph{mudrā}.

Following the \emph{Haṭhapradīpikā}, the next major section of the manuscript bundle comprises Koraṇṭaka's three \emph{paddhati}s on Haṭhayoga, namely, \emph{Koraṇṭakapaddhati~I}, \emph{II} and \emph{III}. These texts are the focus of the present volume.

A half-folio containing what appear to be scribal notes occurs between \emph{Koraṇṭakapaddhati}~\emph{I} and~\emph{II.} In Table~\ref{introtable1}, this half-folio has been designated Fragment 1. As shown in Figure~\ref{introfig3}, Fragment 1 contains two small illustrations in red ink of yogis performing \emph{āsana}s. The accompanying text, written informally in red ink, resembles a scribal note rather than part of the copied manuscript. Unlike the surrounding Sanskrit texts, it is written in Kannada and refers to the seven \emph{āsana}s taught in the \emph{Jñānaratnamālā}, a yoga text located at the end of the manuscript bundle (see Table~\ref{introtable1}). The note written across the top of the folio states, ``Because they are not available in this book, seven \emph{āsana} names have been taken from the \emph{Jñānaratnamālā}.''\endnote{\emph{Mysore Yoga Manuscript}, f. 72\emph{v}: \emph{i grantha ralli illade iruva karanāṃtare jñānaratnamālā tegedukonda} 7 \emph{āsana nāmagalu}. The words \emph{i grantha ralli illade} have been understood as \emph{i grantha dalli illade}. We would like to thank H.\,H. Maharani Pramoda Devi, Prof.~C.\,S. Yogananda and Raghuvendra, M.\,S. for their help with deciphering and translating this statement in Kannada.} Below this note, the following names of seven postures are written across the page: \emph{dhenu/gomukha}, \emph{yogāsana}, \emph{siṃhāsana}, \emph{bhadrāsana}, \emph{vīrāsana}, \emph{padmāsana}, and \emph{siddhāsana.} Illustrations of \emph{siṃhāsana} and \emph{bhadrāsana} have been added, and it appears that space was left for other illustrations that were never completed.\endnote{The transcript of the \emph{āsana} names on f.~72\emph{v} is: 1 \emph{dhenurāsanaṃ gomukhareṃ runā oṃkara} 1 \emph{yogāsana} 1 \emph{siṃhāsana} 1 \emph{bhadrāsana} 1 \emph{vīrāsana} 1 \emph{padmāsana} 1 \emph{siddhāsana}.}




Following \emph{Koraṇṭakapaddhati}~\emph{III}, the bundle contains another textual fragment (designated Fragment~2 in Table~\ref{introtable1}), consisting of a short dialogue between Śiva and the goddess. It begins with the invocation, ``Homage to the glorious Gaṇeśa'' (\emph{śrīgaṇeśāya namaḥ}) and ends with the statement, ``Thus it is complete'' (\emph{iti samāptaṃ}). The text bears no title, and we have not been able to identify it. Its incomplete state suggests that it is a fragment of a larger work.

The next text in the manuscript bundle is the \emph{Deśikastotra}, which is a short hymn venerating Śiva in his form as Śarabha, the guru. The colophon indicates that it is part of the \emph{Ākāśabhairavakalpa},\endnote{\emph{Mysore Yoga Manuscript}, f.~81\emph{v}: \emph{iti śrī ākāśabhairavakalpe sarvasiddhiprade śrīumā\-maheśvarasaṃvade deśikastotraṃ saṃpūrṇaṃ.}} a large tantric compendium on the ritual and healing procedures associated with the fearsome winged deity Ākāśabhairava, who was frequently invoked for exorcism (Goudriaan and Gupta 1981: 115). Although the \emph{Ākāśabhairavakalpa} is cited in the \emph{New Catalogus Catalogorum} as the source of various shorter texts, the \emph{Deśikastotra} preserved in the \emph{Mysore Yoga Manuscript} may derive from this work because it closely resembles a chapter of the same name in a Nepalese manuscript of the \emph{Ākāśabhairavakalpa}.\endnote{Eighty-one chapter names of the \emph{Ākāśabhairavakalpa} in the Nepalese manuscripts A 144/13 and A 145/9 (NGMCP) have been listed on the blog of Michael Slouber: \url{https://garudam.info/akasabhairavakalpa/} (accessed 20 March 2026). We would like to thank Michael Slouber for allowing us to consult his copy of manuscript A 144/13.}

Likewise, the \emph{Yantradvāra,} the penultimate text in the \emph{Mysore Yoga Manuscript}, may also derive from the same tantric compendium, as its colophon attributes it to the \emph{Śrīākāśakalpa.}\endnote{\emph{Mysore Yoga Manuscript} f. 85\emph{r}: \emph{iti śrīākāśakalpe umāmaheśvarasaṃvāde yantradvāra}. There is a chapter called \emph{mantrayantraprakrama} in the Nepalese manuscript of the \emph{Ākāśabhairavakalpa} (A~144/13). As it is not the focus of this book, we have not compared it to the \emph{yantradvāra} passage in the \emph{Mysore Yoga Manuscript}.}

As for the remaining fragments, Fragment 3 repeats the opening portion of the preceding text, the \emph{Deśikastotra}, whereas Fragment 4 is another short, untitled text that begins with the invocation, ``Homage to the glorious Gaṇeśa'' (\emph{śrīgaṇeśāya namaḥ}) and lacks a final colophon. It may preserve a passage from the \emph{Ākāśabhairavakalpa}, as it opens by stating that it concerns the ritual procedures for the worship of Śarabhasālba, more commonly spelt Śarabhasāluva, one of the forms of the deity Ākāśabhairava.\endnote{\emph{Mysore Yoga Manuscript}, f. 86\emph{r}: \emph{atha śarabhasālbapūjau prakaraṇavidhānaṃ}. On the forms of Ākāśabhairava in the \emph{Ākāśabhairavakalpa}, see Goudriaan and Gupta 1981: 115.}

The third-to-last text in the bundle, the \emph{Śivanṛtya}, is a tantric work in nine chapters on the construction and application of ritual diagrams (\emph{yantra}). Numerous manuscripts of this text are reported in the \emph{New Catalogus Catalogorum}, which notes that it is also known as the \emph{Śivatāṇḍava}. The version preserved in the \emph{Mysore Yoga Manuscript} contains ten square numerical grids, ranging from three to nine rows and columns, inserted at various points of the text.

The \emph{Jñānaratnamālā}, the final text in the \emph{Mysore Yoga Manuscript}, is a work of six chapters on Haṭhayoga. Composed in a medieval Hindi-inflected vernacular, it includes sections on the \emph{ṣaṭkarma}, \emph{āsana}, \emph{prāṇāyāma}, \emph{mudrā}, and \emph{bandha}. Illustrations depicting six \emph{cakra}s, four seated yogis, and a yogi performing \emph{mahāmudrā} accompany this text.

It is hoped that the illustrated \emph{Haṭhapradīpikā} and \emph{Jñānaratnamālā} preserved in the \emph{Mysore Yoga Manuscript} will be published as subsequent volumes within this same series.
\vskip 18pt

\Subsection{Numbering of Folios and Illustrations}\label{introsec1.4}
\vskip 9pt

\noindent
An unusual codicological feature of the \emph{Mysore Yoga Manuscript} is its dual foliation, which departs from conventional practice. Rather than numbering the \emph{verso} (back) side of each folio, as is customary in South Asian manuscripts, the \emph{verso}s are unnumbered. Instead, two separate sequences of folio numbers appear in the margins of the \emph{recto} (front) side of each folio.


As shown in Figure~\ref{introfig4}, the first set of folio numbers is written in black ink in the upper left margin of the \emph{recto} side. This set was likely added when the individual texts were originally copied and, as shown in Table~\ref{introtable2}, corresponds to the present arrangement of the folios within the bundle. Thus, the first text, the \emph{Haṭhapradīpikā}, is numbered 1--32, while the second, \emph{Koraṇṭakapaddhati~I}, is numbered 1--38. \emph{Koraṇṭakapaddhati}~\emph{II} and~\emph{III} were originally assigned the folio numbers 1--4, but these were crossed out and replaced with the numbers 39--42. This alteration suggests that a later scribe regarded all three \emph{paddhati}s as a single text, perhaps because the \emph{Koraṇṭakapaddhati~I} lacks a colophon and \emph{Koraṇṭakapaddhati}~\emph{III} ends with two large flourishes (\emph{puṣpikā}) in the left and
%\end{nohyphens}
\end{eng}
\end{multicols*}



\begin{center}
\refstepcounter{figure}
\includegraphics[height=0.97\textheight]{intro/04.jpg}\\
{Figure 4: \emph{Koraṇṭakapaddhati~I}, f.\ 12\emph{v}–13\emph{r}, \emph{āsana} nos.\ 41–44.}
\label{introfig4}
\end{center}

\newpage

\begin{table}[h]
\centering
\renewcommand{\arraystretch}{1.6}
\caption{The Black and Red Folio Numbering of the \emph{Mysuru Yoga Manuscript}}\label{introtable2}
\vskip 6pt

\begin{tabular}{clcl}
\hline
\rowcolor{tableshade}
\multicolumn{1}{p{4cm}}{\centering Black Folio Numbers (current sequence)} & 
\multicolumn{1}{c}{Name of Text} & 
\multicolumn{1}{p{4cm}}{\centering Red Folio Numbers (earlier sequence)} &
\multicolumn{1}{c}{Name of Text} \tabularnewline
\hline
1--32 & \emph{Haṭhapradīpikā} & 1--32 & 1st half of \emph{Koraṇṭakapaddhati} \emph{I} \\
1--38 & \emph{Koraṇṭakapaddhati} \emph{I} & 33--64 & \emph{Haṭhapradīpikā} \\
39--42 & \emph{Koraṇṭakapaddhati} \emph{II} and \emph{III} & 65--70 & 2nd half of \emph{Koraṇṭakapaddhati} \emph{I} \\
& & 71--74 & \emph{Koraṇṭakapaddhati} \emph{II} and \emph{III} \\
\hline
\end{tabular}
\end{table}

\begin{multicols*}{2}
\begin{eng}
\noindent
right margins. The folio numbering in black ink resumes at 1 again for each of the subsequent texts.

The second set of folio numbers is inscribed in the right margin of the \emph{recto} side. Written in red ink and enclosed within circles, these numbers were most likely added when the individual texts were first assembled together with the intention of creating a manuscript bundle. This sequence indicates that the original order of the texts differed slightly from the present arrangement, which follows the sequence recorded by the black folios numbers.

In the present arrangement of the manuscript bundle, the red folio numbers begin at number 33 on the first folio of the \emph{Haṭhapradīpikā} and continue through to number 64, the final folio of the text\emph{.} The first half of \emph{Koraṇṭakapaddhati~I} is numbered 1--32, whereas the remainder of the text is numbered~65--70. This discontinuous numbering reveals that the manuscript bundle originally opened with the first half of \emph{Koraṇṭakapaddhati~I}, which contains the illustrations of one hundred and fourteen \emph{āsana}s and the teachings on \emph{ṣaṭkarma}, followed by the \emph{Haṭhapradīpikā} and then the remainder of \emph{Koraṇṭakapaddhati~I}. Before the manuscript was bound in its present form, however, the order of these two texts was corrected, resulting in the discontinuity in the red foliation. The remaining folios are numbered consecutively from 71 to 91 in red, reflecting the present arrangement of the manuscript.

\emph{Koraṇṭakapaddhati}~\emph{I} bears two additional sets of numbers, which enumerate the descriptions and illustrations of the \emph{āsana}s, respectively. There is, however, a discrepancy between these two sets. The descriptions of the \emph{āsana}s are numbered 1--112, whereas the illustrations are numbered 1--114. There are several errors in the numbering of the \emph{āsana} descriptions: the numbers 52, 86, and 87 are each repeated, while the number 111 is omitted. Consequently, despite the erroneous numbering, \emph{Koraṇṭakapaddhati}~\emph{I} contains a total of 114 descriptions and 114 illustrations of \emph{āsana}s.

In the \emph{Haṭhapradīpikā} and \emph{Koraṇṭakapaddhati}~\emph{I}, a third set of numbers, totalling one hundred and twenty-two, appears in the left and right margins alongside some of the \emph{āsana} illustrations. Written in larger numerals and in a different hand from the other sets of numbers mentioned above, this series corresponds to the \emph{āsana}s selected for reproduction in the \emph{Śrītattvanidhi}. In the \emph{Śrītattvanidhi}, these selected \emph{āsana}s are reproduced according to this numbering sequence and are deliberately presented in a different order from that found in the \emph{Haṭhapradīpikā} and \emph{Koraṇṭakapaddhati}~\emph{I} (see Table 6). One set of these numbers is written in black ink, while the other is in red ink and enclosed within squares. The black numbers correspond to the first eighty \emph{āsana}s selected for inclusion in the \emph{Śrītattvanidhi}, which are designated the primary list. The red numbers enclosed within squares correspond to the remaining forty-two \emph{āsanas} selected for inclusion, which are classified as secondary in the \emph{Śrītattvanidhi}. This complete set of numbers was almost certainly added by the person who compiled the chapter on \emph{yogāsana} in the \emph{Śrītattvanidhi}, which contains descriptions and illustrations of a total of one hundred and twenty-two \emph{āsana}s that were all sourced from the \emph{Mysore Yoga Manuscript.}
\vskip 18pt

\Subsection{The Source of the \emph{Saṅkhyāratnamālā} and \emph{Śrītattvanidhi}}\label{introsec1.5}
\vskip 9pt

\noindent
This section examines the textual relationship between the \emph{Mysore Yoga Manuscript} and two important compilatory works produced at the court of Krishnaraja Wadiyar III: the \emph{Saṅkhyāratnamālā} and the \emph{Śrītattvanidhi.}

The \emph{Saṅkhyāratnamālā} is an illustrated numerical lexicon that was completed on Thursday, 29~March 1849~CE.\endnote{Column 14 of \emph{Descriptive Catalogue of Sanskrit Manuscripts} by Malleddevaru et al.\ (1987, 462--63) states \emph{saumya, caitra, śu 5 guruvāra}, which can be understood as ``Thursday, on the 5th lunar day, in the bright half of the month \emph{caitra} in the year, \emph{saumya} {[}in the 60 year' cycle of Jupiter{]}.'' As far as we know, only one copy of this work was made, so we assume this date refers to the date of its composition. Hāvanūra (2011: 179) also dates the work at 1849~CE.} It is regarded as one of the important works attributed to Krishnaraja Wadiyar~III,\endnote{For example, Hāvanūra 2011: 179, who refers to this work as both the \emph{Saṅkhyāratnamālā} and \emph{Saṅkhyāmāle}. See also Narasimhacharya 1974: 173 and Rao 1922: 172 (where it is called the \emph{Saṅkhyāmāle}).} and contains artwork redolent of that in other illustrated works produced by the court. The \emph{Saṅkhyāratnamālā} defines each number through a set of traditional correspondences. For example, the number `8' is associated with the eightfold auxiliaries (\emph{aṅga}) of yoga (\emph{aṣṭāṅgayoga}). Somewhat unexpectedly, however, the number `80' is defined by a list of \emph{yogāsana}s. This is noteworthy because one might instead expect \emph{āsana}s to correspond to the number `84', the canonical number of postures said to have been first taught by Śiva (Birch 2020: 107--108).

As shown in Table~\ref{introtable5} (Appendix~A), the \emph{Saṅkhyāratnamālā}'s list of \emph{āsana}s is an amalgamation of postures drawn from the \emph{Koraṇṭakapaddhati~I} (KP) and the \emph{Haṭhapradīpikā} (HP) of the \emph{Mysore Yoga Manuscript}. Unlike the \emph{āsana}s in the \emph{Śrītattvanidhi}, however, those in the \emph{Saṅkhyāratnamālā} are not numbered.\endnote{Table~\ref{introtable5} also reproduces the layout of the list of \emph{āsana} names as they appear on the folios of the \emph{Saṅkhyāratnamālā.} On folio 356b, the first forty-four \emph{āsana}s are arranged in two columns, while on folio 357a, thirty are listed in two columns and the final six in a single column.} The heading above the list of postures in the \emph{Saṅkhyāratnamālā} reads, ``The \emph{yogāsana}s in the \emph{Haṭhayogapradīpikā.}'' Before the nineteenth century, the title \emph{Haṭhayogapradīpikā} was rarely used for the text which was more commonly referred to in its colophons as the \emph{Haṭhapradīpikā}. As outlined in Table~\ref{introtable5}, only seven of the \emph{āsana} names in the \emph{Saṅkhyāratnamālā} list derive from the \emph{Haṭhapradīpikā} of the \emph{Mysore Yoga Manuscript}, whereas sixty-seven originate from \emph{Koraṇṭakapaddhati~I}. The remaining six of the eighty \emph{āsana}s may have been drawn from either text.

The inclusion of the postures \emph{cakrāsana} and \emph{yogāsana} in the \emph{Saṅkhyāratnamālā}'s list is particularly telling, as both are absent from standard recensions of the \emph{Haṭhapradīpikā} and do not occur in Koraṇṭaka's \emph{paddhati}s. Both of these \emph{āsana}s, however, are found in the unique version of the \emph{Haṭhapradīpikā} preserved in the \emph{Mysore Yoga Manuscript}. This provides evidence that the Mysore version of the \emph{Haṭhapradīpikā} served as the source for the seven postures. The compiler of the \emph{Saṅkhyāratnamālā}'s list, however, may have attributed the entire list of eighty \emph{āsana}s to the ``\emph{Haṭhayogapradīpikā}'' because the \emph{Haṭhapradīpikā} appears as the first text in the \emph{Mysore Yoga Manuscript}.

Table~\ref{introtable5} further shows that every \emph{āsana} listed in the \emph{Saṅkhyāratnamālā} also occurs in the \emph{Śrītattvanidhi} (ff.\ 834--851)\emph{.} An even closer relationship exists between the \emph{Mysore Yoga Manuscript} and the \emph{Śrītattvanidhi}.

Composed between 1810 and 1868~CE, the \emph{Śrītattvanidhi} is a textual and pictorial compendium of divine iconography and iconometry, illustrated with painted miniatures and drawings. It provides technical instructions for artists on depicting gods, goddesses, and mythological figures, and is divided into nine \emph{nidhi}s (``treasures'').\endnote{The nine \emph{nidhi}s are: (1) \emph{Śaktinidhi}, (2) \emph{Viṣṇunidhi}, (3) \emph{Śivanidhi}, (4) \emph{Brahmanidhi}, (5)~\emph{Grahanidhi}, (6) \emph{Vaiṣṇavanidhi}, (7) \emph{Śaivanidhi}, (8) \emph{Āgamanidhi}, and (9) \emph{Kautukanidhi}.} The seventh \emph{nidhi}, the \emph{Śaivanidhi}, contains a section on \emph{yogāsana}s, whose textual descriptions are identical to the corresponding descriptions in the \emph{Mysore Yoga Manuscript}. It is possible that the \emph{Śaivanidhi} postdates the \emph{Saṅkhyāratnamālā}, because the latter is reportedly cited by name in the \emph{Grahanidhi}, which is the earlier fifth volume of the \emph{Śrītattvanidhi}.\endnote{This information was provided by the Mysore Oriental Research Institute, but we have not been able to verify this ourselves by viewing the manuscript.}

As shown in Table~\ref{introtable6} (Appendix A), the one hundred and twenty-two \emph{āsana}s of the \emph{Śrītattvanidhi} are likewise an amalgamation of postures drawn from the \emph{Koraṇṭakapaddhati}~\emph{I} (KP) and the \emph{Haṭhapradīpikā} (HP) of the \emph{Mysore Yoga Manuscript}. As discussed in the previous section, the third set of scribal numbers added to the \emph{āsana} illustrations in the \emph{Mysore Yoga Manuscript} affirms the direction of borrowing. These numbers correspond to the order and numbering of the \emph{āsana}s in the \emph{Śrītattvanidhi} and were almost certainly added when its chapter on \emph{yogāsana}s was compiled\emph{.}

The redactor of the \emph{Śrītattvanidhi} rearranged Koraṇṭaka's \emph{āsana}s, dividing them into two categories, designated primary and secondary, and supplemented them with ten \emph{āsana}s from the \emph{Haṭhapradīpikā}. In the \emph{Mysore Yoga Manuscript}, the redactor's selection of the primary \emph{āsana}s and their order is indicated by numbers in black ink and the secondary ones are marked by numbers in red ink encased in circles.\endnote{\emph{Śrītattvanidhi} (\emph{Śaivanidhi}): ``The principal {[}postures{]} numbering eighty are \emph{āsana}s suitable for yoga. After these, there are also Vṛṣapādakṣepāsana and Trikuṭīnāsana, as well as others; thus, these forty-two \emph{āsana}s are indeed taught here as additional'' (\ldots{} \emph{mukhyāny aśītisaṅkhyāni yogayogyāsanāni hi} \textbar\textbar{} \emph{ataḥ paraṃ cāpi vṛṣapādakṣepāsanaṃ tathā \textbar{}} {[}\ldots{]} \textbar\textbar{} \emph{trikuṭīnāsanaṃ cānyād ity evam adhikāny api} \textbar\textbar{} \emph{pradarśitāny āsanāni dvicatvāriṃśad atra vai}).}

It is noteworthy that all but one of the eighty principal \emph{āsana}s in the \emph{Śrītattvanidhi} correspond to those listed in the \emph{Saṅkhyāratnamālā.} The sole exception is that the Pose of Viṣṇu's Three Steps (\emph{trivikramāsana}), which replaces the Bow Pose (\emph{dhanurāsana}) among the principal \emph{āsana}s, while the Bow Pose (\emph{dhanurāsana}) is assigned to the secondary list in the \emph{Śrītattvanidhi.} The reasons for this substitution are unknown.

As outlined in Table~\ref{introtable6}, one hundred and twelve of the \emph{āsana}s in the \emph{Śrītattvanidhi} were drawn from \emph{Koraṇṭakapaddhati~I} and ten from the \emph{Haṭhapradīpikā}. In cases where an \emph{āsana} occurs in both source texts, such as the Upturned Turtle Pose (\emph{uttānakūrmāsana}), the scribal numbering in the margins of the \emph{Mysore Yoga Manuscript} makes it possible to identify with confidence the source from which it was copied. The process of compilation is also evident in the illustrations of the \emph{āsana}s in the \emph{Śrītattvanidhi}, which combines elements of two distinct artistic styles. The illustrations derived from the \emph{Haṭhapradīpikā} often depict the yogi with long, loose hair rather than the distinctive tight topknot, which is characteristic of the illustrations of \emph{Koraṇṭakapaddhati~I}.

In compiling the \emph{Śrītattvanidhi}'s chapter on \emph{yogāsana}, the redactor did not preserve the sequence of the postures presented in \emph{Koraṇṭakapaddhati~I}. The rationale for this reordering remains unclear. The introduction to this \emph{yogāsana} chapter of the \emph{Śrītattvanidhi} states that the eighty principal postures are found exclusively in yoga scripture.\endnote{\emph{Śrītattvanidhi} (\emph{Śaivanidhi}): ``Yoga postures that are exclusive to scripture on yoga should be known as eighty. Now, the shapes of the eighty postures are written down in the manner of a yoga scripture'' (\emph{yogāsanaṃ yogaśāstramātraṃ jñeyam aśītidhā \textbar\textbar{} atha aśītyāsanasvarūpāṇi yogaśāstrarītyā likhyante}). See the previous footnote for the comment on the suitability of the principal \emph{āsana}s for yoga.} At first glance, this claim appears to contradict the fact that both the principal and secondary postures were drawn from two Haṭhayoga texts. One possibility is that the redactor was aware of another yoga scripture, now lost or otherwise unknown, that defined or categorised a corpus of eighty principal \emph{āsana}s. Whatever the case, the evidence suggests that the redactor attached little importance to preserving the original order of the postures in \emph{Koraṇṭakapaddhati~I}. Instead, the collection was deliberately reordered and augmented with ten \emph{āsana}s from the \emph{Haṭhapradīpikā}, seven of which were incorporated into the category of primary postures.
\end{eng}
\end{multicols*}
\newpage
~
\newpage
\null\vskip 0.1cm

\begin{eng}
\centerline{\Section{The Three Yoga Manuals of Koraṇṭaka: \emph{Koraṇṭakapaddhati~I--III}}\label{introsec2}}
\end{eng}
\vskip 22pt

\begin{multicols*}{2}
\begin{eng}
\Subsection{Who was Koraṇṭaka?}\label{introsec2.1}
\vskip 9pt

\lettrine[lines=3,lhang=0.04,nindent=0em]{T}{he three manuals} in this volume attribute their teachings to a perfected yogi (\emph{siddha}) who is referred to by the names, Kapālakuraṇṭaka, Koraṇṭaka, Goraṇṭaka, and Kuraṇṭaka. While the first of these names is distinguished by the word \emph{Kapāla}, the latter three are likely orthographic variants shaped by vernacular usage. The evidence discussed below indicates that these names refer to the same \emph{siddha}, who later became a prominent figure in the Maharashtrian Nātha tradition. Since the Sanskrit form ``Koraṇṭaka'' most closely corresponds to the Marathi and Kannada spellings, it has been adopted here to refer to the \emph{siddha} associated with these three manuals, which were most likely composed for a Vaiṣṇava ascetic order in eighteenth-century Maharashtra.

The earliest references to this \emph{siddha} occur in Marathi, Kannada, Telugu, and Sanskrit sources, where his name appears in several orthographic forms. In premodern Marathi literature, he is called Koraṇṭaku, a spelling found in the thirteenth-century \emph{Līḷācaritra} (528) and fourteenth-century \emph{Tattvasāra} (872).\endnote{We wish to thank Amol Bankar for the reference in the \emph{Līḷācaritra} and \emph{Tattvasāra}.} In Kannada literature, the thirteenth-century \emph{Siddharāmacāritra~}(1.10,~9.43) refers to him as Kōrāṇṭa and Kōrāṇṭaka, while the sixteenth-century \emph{Revaṇasiddheśvarapurāṇa} (6.64) calls him Kōrāṇṭaka\endnote{We wish to thank James Mallinson and Nishith Amarnath for the reference in the \emph{Siddharāmacāritra}.} In the fourteenth-century Telugu work, the \emph{Navanāthacaritramu}, he appears as Goraṇṭakuḍu, and the seventeenth-century Sanskrit yoga text, the \emph{Haṭharatnāvalī}, which was probably composed in Andhra Pradesh (Reddy~1982:~14--17), preserves the forms Gorandaka or Korandaka in its manuscript transmission.

Although these sources differ in their spelling of his name, they consistently associate the \emph{siddha} with Maharashtra. Premodern Marathi and Telugu sources identify him as a disciple of Gorakṣanātha. In the \emph{Navanāthacaritramu} (Jones 2018: 197--8), the story of Goraṇṭakuḍu takes place in the city of Kandāra in Maharashtra, while in the \emph{Līḷācaritra~}(528) he is described as travelling through the sky accompanied by three other \emph{siddha}s, Nāgārjuna, Sevāḷabhakṣu, and Ratnaghoṣā, to the city of Kolhapur, which is in present-day Maharashtra. Taken together, these sources strongly suggest a Maharashtrian provenance, which accords with the transmission of these manuals in Maharashtra and Karnataka (see Section~\ref{introsec2.3}).

His alternative name, Kapālakuraṇṭaka, however, raises the question of Koraṇṭaka's sectarian affiliation. At first sight, this name (literally, the ``skull-bearing Kuraṇṭaka''), together with the prominence of \emph{vajrolīmudrā} in his teachings, a practice associated with the Kāpālikas,\endnote{For example, the \emph{Haṭhapradīpikā}~(3.90) attributes the practice of ``drinking one's urine'' (\emph{amarolī}) to the Kāpālikas of the Khaṇḍa school. \emph{Amarolī,} which is taught in the second manual, was generally regarded as a variation of \emph{vajrolīmudrā}, a sexual practice that is the main technique of the second and third manuals, and is taught at length in the first.} suggests that he was affiliated with a Kāpālika sect. The Kāpālikas were a tradition of Śaiva ascetics for whom the observance of the ``skull vow'' (\emph{kapālavrata}) was emblematic. They were distinguished by carrying a skull bowl and living in cremation grounds, where they performed antinomian rites. That some Kāpālikas may have practised Haṭhayoga is implied in the thirteenth-century \emph{Dattātreyayogaśāstra} (41), while the names Kapālī and Khaṇḍakāpālika in the lineage of \emph{siddha}s at the beginning of the fifteenth-century \emph{Haṭhapradīpikā}~(1.7--8) further suggest an association between Kāpālikas and Haṭhayoga.

Nonetheless, the evidence is insufficient to establish that Koraṇṭaka was himself a Kāpālika. Kāpālikas are known for their worship of Bhairava, a fierce manifestation of Śiva, whereas Koraṇṭaka's manuals appear to have been composed for a Vaiṣṇava ascetic order, since the practice involves the use of the twelve-syllable \emph{mantra} of Vāsudeva and Vaiṣṇava deities (\emph{mūrti}). Moreover, the description of the yogi in the third manual lacks the key insignia of a Kāpālika, such as the skull bowl (\emph{karpara}) and skull staff~(\emph{khaṭvāṅga}).

Rather than indicating sectarian affiliation, the name Kapālakuraṇṭaka may instead be understood in light of the Marathi term \emph{kapāḷakaraṇṭā}, meaning an ``ill-fated person.'' Here, the usage of \emph{kapāla} may derive from the belief that one's fate is written on one's skull or forehead.\endnote{In Sanskrit literature, this is usually referred to by expressions such as \emph{lalāṭalikhita} (``written on one's forehead''). We wish to thank Shaman Hatley for alerting us to this (p.c. 24.7.2026).} The epithet may therefore allude to Koraṇṭaka's untimely death, a theme that appears in later narrative traditions.\endnote{We wish to thank Amol Bankar for informing us of the meaning of the Marathi term \emph{kapāḷakaraṇṭā} (p.c. March 2026).} Thus, in the \emph{Navanāthacaritramu} (Jones 2018: 197--8), he was interred alive in a temple, sacrificing himself to restore the kingdom of Kṛṣṇakandāra, and is never seen again. Similarly, in the \emph{Siddharāmacāritra}, a proto-Vīraśaiva work, he falls into a lake and drowns after failing to worship the \emph{liṅga}.
\vskip 18pt

\Subsection{The Name of Koraṇṭaka's Texts}\label{introsec2.2}
\vskip 9pt

\noindent
On the basis of the available colophons and the likely provenance of both the \emph{siddha} and the \emph{paddhati}s (see Section~\ref{introsec2.3}), the three manuals in the \emph{Mysore Yoga Manuscript} are referred to in this volume as \emph{Koraṇṭakapaddhati~I}, \emph{II}, and~\emph{III}.

Only \emph{Koraṇṭakapaddhati}~\emph{II} and \emph{III} conclude with colophons, which refer to the manuals as \emph{Koraṇṭakapaddhati} and \emph{Kuraṇṭakapaddhati}, respectively.\endnote{The colophon of the second \emph{paddhati}, ``Thus, \emph{vajrolīmudrā} in Koraṇṭaka's manual'' (\emph{koraṇṭakapaddhatau vajrolīmudrā}) suggests that this text is a section of the \emph{Koraṇṭakapaddhati}.} The first \emph{paddhati} lacks a colophon and ends abruptly after its explanation of the tenth \emph{mudrā}, namely \emph{viparītakaraṇī}. This suggests that the manual is incomplete, particularly given that there is no conclusive discussion on meditative absorption (often referred to as \emph{rājayoga} or \emph{samādhi}). Despite the absence of a colophon, a marginal note in a different scribal hand on the first folio identifies \emph{paddhati} \emph{I} as both the \emph{Vajrolīyogapaddhati} and \emph{Kuraṇṭakapaddhati}. By contrast, the opening sentence of \emph{paddhati} \emph{I} states, ``This is a manual on the practice of Haṭha (\emph{haṭhābhyāsapaddhati}) composed by Kapālakuraṇṭaka.'' In light of the marginal note, the expression \emph{haṭhābhyāsapaddhati} appears to describe the nature of the work rather than its title. Therefore, the available colophons and marginal note indicate that all three manuals circulated, at the very least in Mysore, under the titles \emph{Koraṇṭakapaddhati} or \emph{Kuraṇṭakapaddhati.}
\vskip 18pt

\Subsection{Provenance}\label{introsec2.3}
\vskip 9pt

\noindent
It is rare for premodern works on yoga to be attributed to a named author and to make explicit references to geographic regions, landmarks, and historical events. Nevertheless, there are several indications that the three \emph{Koraṇṭakapaddhati}s originated in Maharashtra. At the most general level, all three are attributed to Koraṇṭaka, whose provenance appears to be Maharashtrian; in addition, two of the \emph{paddhati}s contain internal evidence pointing to a Maharashtrian setting.

The provenance of \emph{Koraṇṭakapaddhati}~\emph{I} can be inferred from a statement on the name of the grass used as a probe in the practice of \emph{vajrolīmudrā}. Although writing in Sanskrit, the author gives its Marathi name, \emph{lavālā}, and states that it was well known in Maharashtra and elsewhere.\endnote{\emph{Koraṇṭakapaddhati~I}: ``Resembling a sprout of the jasmine plant, {[}the grass{]} named \emph{haritaśara} is known in Maharashtra and elsewhere as \emph{lavālā}'' (\emph{jātyaṅkurasadṛśo haritaśaraḥ nāma lavālā iti mahārāṣṭrādau prasiddhaḥ}).} This comment suggests that the author was familiar with both Marathi and the use of this grass in the Maratha Empire, whose rule in the eighteenth century extended beyond the borders of present-day Maharashtra.\endnote{On the predominance of the Maratha empire in the eighteenth century, the so-called second Maratha empire that emerged out of regional resistance to the late Mughal occupation, see Guha~2019.}

A reference in \emph{Koraṇṭakapaddhati}~\emph{II} indicates that it too originates from Maharashtra. It mentions Gahananātha, also known as Gahinīnātha and Gāyaṇinātha, who is regarded as the guru of Jñāneśvara's guru. Jñāneśvara was a thirteenth-century poet-saint who composed the Marathi Jñāneśvarī and is considered one of the founders of the Maharashtrian Bhāgavata (Vārkarī) movement (Kiehnle~2005). Gahananātha, therefore, appears to have belonged to the Maharashtrian lineage of the Nātha tradition. As for the third \emph{paddhati}, no comparable evidence is available for its provenance, apart from Koraṇṭaka's broader association with Maharashtra and Karnataka.

The three \emph{paddhati}s attributed to Koraṇṭaka were most probably composed in the eighteenth century. It is certain that \emph{Koraṇṭakapaddhati}s~\emph{I} and \emph{II} postdate the fifteenth century, as the first borrows a verse from the \emph{Haṭhapradīpikā} and part of the second's description of the water enema is calqued on the corresponding passage in the \emph{Haṭhapradīpikā}.\endnote{The verse on the eight \emph{kumbhaka}s in \emph{Koraṇṭakapaddhati}~\emph{I} = \emph{Haṭhapradīpikā} 2.44. \emph{Koraṇṭakapaddhati}~\emph{II} 13ab \textasciitilde{} \emph{Haṭhapradīpikā} 2.26ab.} The \emph{terminus ante quem} for all three manuals is provided by the \emph{Mysore Yoga Manuscript}, which was likely produced during the initial phase of Krishnaraja Wadiyar III's literary programme in the early nineteenth century (see Section~\ref{introsec1.1}). All three \emph{paddhati}s exhibit sufficient signs of textual deterioration to suggest a relatively long transmission,\endnote{Specific evidence of scribal errors, interpolations, omissions, corrupted readings, etc., have been noted in the endnotes to the Sanskrit edition of all three manuals.} indicating that these works had been copied repeatedly for some time prior to their inclusion in \emph{Mysore~Yoga~Manuscript}. As far as we are aware, none of the three \emph{paddhati}s has been cited in any compendium or work on yoga composed before the nineteenth century. It is also possible, albeit less likely, that they were written earlier than the eighteenth century but had a limited circulation.

Given their shared tradition and content (Section~\ref{introsec2.4}), it is likely that all three \emph{paddhati}s originated within the same yogic milieu.
\vskip 18pt

\Subsection{Salient and Unique Features}\label{introsec2.4}
\vskip 9pt

\noindent
Beyond their association with the \emph{siddha} Koraṇṭaka, the three \emph{paddhati}s are connected by shared teachings, which are sometimes articulated differently in each work. Their principal shared teaching is \emph{vajrolimudrā} and the practices associated with it. All three works also prescribe the recitation of the twelve-syllable Vāsudeva \emph{mantra}, which the yogi is instructed to repeat during various cleansing and preparatory practices, including bathing, eating, sleeping, and excreting. Nevertheless, one broad distinction should be noted. \emph{Koraṇṭakapaddhati}~\emph{I} differs considerably from the second and third in terms of the style of composition and structure. \emph{Koraṇṭakapaddhati}~\emph{I} was composed in prose and organised as a set of auxiliaries, with \emph{vajrolī} included as one of ten seals (\emph{mudrā}), whereas \emph{Koraṇṭakapaddhati}~\emph{II} and \emph{III} were composed in verse and mention only one seal, \emph{vajrolī.} The discussion below first compares \emph{Koraṇṭakapaddhati}~\emph{I} and \emph{II}, whose relationship is relatively clear, before turning to the more distinctive features of \emph{Koraṇṭakapaddhati~III}.
\columnbreak

\SSubsection*{\emph{Koraṇṭakapaddhati~I}}\label{koraux1e47ux1e6dakapaddhati-i}
\vskip 7pt

\noindent
\emph{Koraṇṭakapaddhati~I} is the largest of the three manuals. It is organised around a set of auxiliaries, much like an account of \emph{ṣaḍaṅga} or \emph{aṣṭāṅgayoga.} Its contents include the yogi's huts (\emph{maṭhikā}), behavioural codes~(\emph{yama}) and observances (\emph{niyama}), yogic postures (\emph{āsana}), cleansing techniques (\emph{ṣaṭkarma}), breath control (\emph{prāṇāyāma}), and the ten seals (\emph{mudrā}). Among these topics, the discussions of yogic posture and the seal called \emph{vajrolī} are the most extensive. The section on yogic posture constitutes nearly forty percent of the entire work, while the description of \emph{vajrolī} is longer than the combined discussions of all the other seals.

\emph{Koraṇṭakapaddhati~I} contains descriptions of previously undocumented practices, as well as many new details of techniques already known from earlier Haṭhayoga texts. Whereas earlier works such as the \emph{Dattātreyayogaśāstra} (54--57), \emph{Yogayājñavalkya} (5.6--8) and \emph{Haṭhapradīpikā} (1.12--13) describe only one small hut for the practice of yoga, the opening section of \emph{Koraṇṭakapaddhati~I} describes a series of six huts (\emph{maṭhikā}), made of various materials and intended for different purposes. These include a hut approximately five and a half metres high for practising rope poses such as Climbing Up to Heaven (\emph{svargāsana}, no.~99).

The section on behavioural guidelines (\emph{yama}) and observances (\emph{niyama}) lists twenty-five such precepts from the \emph{Bhāgavatapurāṇa}, and then provides the only known commentary on \emph{yama} and \emph{niyama} written specifically for practitioners of Haṭhayoga (see also Section~\ref{introsec2.5}).

\emph{Koraṇṭakapaddhati~I} describes the most extensive and sophisticated practice of complex, moving \emph{āsana}s of all the available works on yoga composed before the twentieth century. Its \emph{āsana}s are categorised under six headings, namely Supine (\emph{uttāna}), Prone (\emph{nyubja}), Stationary (\emph{sthāna}), Standing (\emph{utthāna}), Postures with Ropes (\emph{rajju}), and Postures that Pierce the Sun and Moon (\emph{sūryacandrabhedana}) (see Section~\ref{introsec3.2}, Table~\ref{introtable3}). The names of several of these categories appear to have been derived from postures that recur frequently within them. In the Supine group (\emph{uttānāsana}), for example, a supine position links thirteen of the twenty-two \emph{āsana}s, each of the thirteen beginning with the statement, ``Having lain supine'' (\emph{uttānaṃ śayanaṃ kṛtvā}). This group also includes postures with rotating, rolling, standing, inverting, back-bending and tumbling movements. It therefore appears that the term ``supine'' refers not to a strict category of similarly shaped postures, but rather to the reoccurring transitional position in the sequence. This also appears to be true of the Prone (\emph{nyubja}) group. An interesting exception is the rope category, in which all the postures require a rope, although one is also performed with a stone weight.

Although referred to as the six cleansing practices (\emph{ṣaṭkarma}), \emph{Koraṇṭakapaddhati~I} actually describes nine, including two that are not mentioned in the \emph{Haṭhapradīpikā}, which is the earliest known work to codify the \emph{ṣaṭkarma}. These are \emph{bhrāmaṇakriyā} (also called \emph{ādhāraśuddhikriyā} and \emph{gaṇeśakriyā}), which aims to clean the rectum and anus, and \emph{manthanapraveśa}, which cleanses the urethra by inserting and churning various probes. Unlike other yoga texts, \emph{Koraṇṭakapaddhati~I} claims that \emph{nauli} can be perfected by repeatedly clenching and relaxing the anus like a horse and that \emph{trāṭaka} can be performed on various gazing points to obtain specific benefits, such as determining whether gems are fake or genuine.

Koraṇṭaka's advice on diet is unique among yoga traditions and appears to aim at cultivating indifference toward food, much like \emph{vajrolimudrā} is aimed at cultivating indifference toward women:

\begin{quote}
Food should be eaten very quickly. {[}The yogi{]} should not take pleasure in food and so forth, and he should not take what is harmful. He should eat food as though it were medicine.\endnote{For this passage in the edition, see page~\pageref{food}.}
\end{quote}

\emph{Koraṇṭakapaddhati~I} teaches eight breath retentions (\emph{kumbhaka}), which largely correspond to those taught in the \emph{Haṭhapradīpikā}. The distinctive features include the incremental breathing of \emph{sūryabhedana,} the silent inhalation and exhalation of \emph{ujjāyī}, the suppression of yawning in \emph{sītkāra,} the omission of rapid breathing in \emph{bhastrikā}, and the manipulation of semen (\emph{bindu}) in \emph{mūrcchā}. Most remarkable is the unique method of moving all the winds that is performed to prepare the yogi for these retentions. This method combines the pelvic, abdominal, and throat locks, known as \emph{mūla}, \emph{uḍḍīyāna}, and \emph{jālandharabandha,} respectively\emph{,} with the forward-bending posture \emph{paścimatānāsana} as well as breath retention and belching.

The ten \emph{mudrā}s of \emph{Koraṇṭakapaddhati~I} are the same as those in the \emph{Haṭhapradīpikā}, although there are various differences, such as the omission of tapping the buttocks on the ground in \emph{mahāvedha}, and the omission of Bellows breath in the practice of \emph{śakticālana}. The practice of the three locks (\emph{bandha}) is also unique in that each involves moving the breath through different regions of the body before performing \emph{prāṇāyāma} (i.e., holding the breath). The notion of moving the breath internally through the practice of \emph{mudrā}s is extended to \emph{viparītakaraṇī}, in which air is expelled through the anus while the body is inverted.

The account of \emph{vajrolimudrā} is the most extraordinary of its kind in any premodern yoga text. It describes a series of preparatory practices for this \emph{mudrā}, beginning with the insertion of stalks of distinct plants and probes made of various substances into the urethra, as deep as ten finger-breadths and for periods up to three hours. These practices cause the yogi acute discomfort, disorientation, and a sharp pain in the bladder. To endure this, the yogi is advised to sing the names of god. He develops the capacity to stop the downward flow of semen through repeatedly stopping and releasing the flow of urine and faeces when excreting. He is also instructed to engage in a series of increasingly intense acts of amorous play with women. The main aim of this complex series of techniques is to cultivate detachment (\emph{vairāgya}) toward women and to maintain celibacy through semen retention rather than through abstinence from sex.
\columnbreak

\SSubsection*{\emph{Koraṇṭakapaddhati~II}}\label{koraux1e47ux1e6dakapaddhati-ii}
\vskip 7pt

\noindent
\emph{Koraṇṭakapaddhati}~\emph{II} focuses on \emph{vajrolīmudrā} and the complex repertoire of practice surrounding it. Firstly, the author explains the preparatory roles of yogic posture, cleansing the anus, the water enema, and purifying the penis. Then follow the practices of amorous play, herbal urethral enemas, \emph{mantra} recitation, the techniques of \emph{sahajolī} and \emph{amarolī}, intercourse with sixteen women, and the Pigeon Throat, Frog Belly, and Cup techniques. Apart from \emph{sahajolī,} \emph{amarolī}, and the Cup technique, all of these topics are taught in \emph{paddhati I}. In effect, \emph{Koraṇṭakapaddhati}~\emph{II} reconfigures and condenses the content of the first \emph{paddhati}, making it clear, for example, that in Koraṇṭaka's tradition of yoga, its dynamic and demanding \emph{āsana} practice served as preparation for \emph{vajrolīmudrā}.

Like \emph{paddhati} \emph{I}, \emph{Koraṇṭakapaddhati}~\emph{II} was clearly intended for a Vaiṣṇava audience. It states that \emph{vajrolī} belongs to Viṣṇu and that the worship of Vāsudeva eliminates all obstacles. The authors of \emph{Koraṇṭakapaddhati}s~\emph{I} and \emph{II} drew upon the same sources, namely, the \emph{Bhāgavatapurāṇa} and \emph{Haṭhapradīpikā}.\endnote{\emph{Paddhati~I} borrows a verse from the \emph{Bhāgavatapurāṇa} on the \emph{yama} and \emph{niyama}, and its verse on the eight \emph{kumbhaka}s is very close to \emph{Haṭhapradīpikā} 2.44. The first verse of \emph{Paddhati} \emph{II} is based on \emph{Bhāgavatapurāṇa} 11.11.14, and its description of \emph{basti} is close to \emph{Haṭhapradīpikā} 2.26ab.} Both works also share a number of practices and distinctive details that appear to be unique to Koraṇṭaka's tradition of yoga. These include the Pawing the Leg like a Bull (\emph{vṛṣapāda}) and the Boar (\emph{varāha}) poses; cleansing the anus; contracting and relaxing the anus like a horse during defecation; inserting probes, tubes, and stalks into the urethra; sleeping on a bed of ashes; amorous play with the same progressive stages (reciting poetry, standing near a woman, touching her, and so forth); the use of herbal enemas; counting repetitions of practice with the twenty-four names of Vāsudeva; and the Pigeon Throat technique, in which the bladder puffs up like a pigeon's throat.

Beyond the material shared with \emph{paddhati} \emph{I}, \emph{Koraṇṭakapaddhati}~\emph{II} also preserves numerous teachings unattested in any other surviving source on yoga. Its unique material encompasses explanations of terminology, attributions of techniques, practical instructions, and accounts of the purported effects of the practice and its soteriological benefits. It explains the derivation of the name \emph{vajrolī}; attributes the revelation of this \emph{mudrā} to Matsyendranātha; states that an adept can perform sixteen thousand finger rotations to cleanse the anus in forty-eight minutes; notes that Kanthaḍī was renowned for teaching the water enema; prescribes the Boar Pose for an anal variation of Bellows Breath; and instructs the yogi to draw mercury through the urethral tube. It also includes tongue sex (\emph{jihvāmaithuna}) and an eightfold sexual intercourse without the final stage of emission; a detailed description of the yogi's consort; the practices of \emph{sahajolī} and \emph{amarolī}; the soteriological benefits of \emph{vajrolī,} such as knowledge of the Self and the conquest of death; the initial side effects of the practice, such as pain, friction in the inner channels, and disturbance of the \emph{cakra}s; the transformation of pain into itching and ultimately joy; the practice of drawing a gem into and expelling it from the urethral tube with different sounds; and the Frog-Belly and Cup techniques.

The final colophon of \emph{Koraṇṭakapaddhati}~\emph{II} suggests that it may originally have formed a section on \emph{vajrolī} within a larger work entitled the \emph{Koraṇṭakapaddhati}. Even so, \emph{Koraṇṭakapaddhati}~\emph{II} can be read as an independent manual, with one possible exception: verse 20 implies that remedies for fever had been explained elsewhere. Nonetheless, since scribal colophons may have been added during the course of a text's transmission, it remains possible that \emph{Koraṇṭakapaddhati}~\emph{II} was originally composed as a work in its own right devoted to \emph{vajrolīmudrā.}
\vskip 14pt

\SSubsection*{\emph{Koraṇṭakapaddhati III}}\label{koraux1e47ux1e6dakapaddhati-iii}
\vskip 7pt

\noindent
Whereas the second manual shares and supplements material found in the first, \emph{Koraṇṭakapaddhati~III} preserves a more distinctive set of teachings that illuminate the social identity and ritual practices of Koraṇṭaka's yogic tradition. It nevertheless features \emph{vajrolīmudrā} and the Vāsudeva \emph{mantra}. \emph{Koraṇṭakapaddhati}~\emph{III} opens with a general discussion of the practice of \emph{āsana} and \emph{prāṇāyāma}, followed by verses describing their health benefits, herbal remedies, and dietary prescriptions. Verse 12 indicates that these practices are preparatory, as they are intended to benefit one engaged in \emph{vajrolī}. The middle section of the third \emph{paddhati} contains what is perhaps the most detailed portrayal of a yogi in any surviving Sanskrit work, describing his dress and accoutrements and thereby providing a rare insight into the material culture and social identity of this tradition. The final section concerns \emph{mantra} recitation using a golden sacrificial brick, a ritual practice unattested in other premodern yoga texts and one that points to a distinctive ritual dimension of Koraṇṭaka's tradition.

\emph{Koraṇṭakapaddhati}s~\emph{I} and \emph{III} share similar dietary prescriptions, and both refer to a person by the name of Kāki, who, to the best of our knowledge, is unattested in other literature. In \emph{Koraṇṭakapaddhati}~\emph{I}, the name Kāki appears in a list of \emph{siddha}s renowned for their scriptures, whereas in \emph{Koraṇṭakapaddhati~III} he is credited with prescribing the Black Oil plant for wind ailments. One possibility is that Kāki is a shortened form of Kākacaṇḍīśvara, a \emph{siddha} whose name appears in the lineage at the beginning of the \emph{Haṭhapradīpikā} (1.7) and in the alchemical work, the \emph{Rasaratnasamuccaya~}(1.6). Kākacaṇḍīśvara is associated with various herbal compounds compiled in the \emph{Kākacaṇḍīśvarakalpatantra}. Alternatively, the name Kāki might refer to a \emph{siddha} who was associated with crows (\emph{kāka}), an association that might explain the sculpture of a \emph{siddha} seated with a bird perched above his left arm in Cave~29 (Gaur Lene, RWD2) of Panhale Kazi in Maharashtra.\endnote{We wish to thank Amol Bankar for his insightful comments on the possible identity of Kāki and the reference to the sculpture at Panhale Kazi.}

Unique material in \emph{Koraṇṭakapaddhati}~\emph{III} includes the claim that the practice of \emph{āsana} alleviates pain and cures imbalances caused by the practice itself. Yogis are also said to be able to cure their illnesses without the assistance of a doctor. When sweat arises from \emph{prāṇāyāma}, the yogi is advised to rub it with ashes, which is slightly different from the injunction in the \emph{Haṭhapradīpikā} (2.13) to rub the limbs with the sweat.\endnote{Cf.\ \emph{Koraṇṭakapaddhati}~\emph{I}, which states that ejaculated semen and sweat should be rubbed into the body during the practice of \emph{vajrolī}.} For the sake of physical wellbeing, the yogi is instructed to wrap himself in a piece of fine cloth and sleep on a bed of ashes. When trembling caused by \emph{prāṇāyāma} subsides, the yogi's channels are said to become pure. A distinct set of benefits of \emph{prāṇāyāma} follows in verses 8 and 9, such as the breath moving lightly in the body and the teeth becoming strong. The exceptions are purity of the eyes and thinness of the body, which are benefits mentioned in the \emph{Haṭhapradīpikā}~(2.80).

The most distinctive feature of \emph{Koraṇṭakapaddhati}~\emph{III}, however, is its elaborate description of the yogi and his ritual practice. The account of the yogi and his accoutrements in \emph{Koraṇṭakapaddhati~III} is remarkable for several reasons. Such descriptions are rare in works on Haṭhayoga, perhaps because these texts were intended for a wide range of practitioners whose appearance and dress could not be described in any definitive way. By contrast, descriptions of the yogi and his insignia occur in sectarian contexts in other literature, such as the Avadhi romances known as the \emph{Premakathā}s (Mallinson 2021). The opening statement in the third \emph{paddhati}'s description of the yogi expresses an indifference towards sectarian markers, stating that marks and emblems should be worn mockingly and playfully. A similar, albeit more straightforward, sentiment is found in the \emph{Amanaska} (2.34--35), a twelfth-century work on Rājayoga, where it is stated that adhering to sectarian signs (\emph{liṅga}) serves no useful purpose.

{\parfillskip=0pt
Many of the accoutrements mentioned in the description of the yogi are emblematic of a wide range of Indian ascetics, such as his use of the black antelope skin, a tiger-skin garment, water pot and staff, and his wearing of matted locks, white ash, a loin cloth, and sandals. However, the interesting details added to these items include black antelope skins that have the horns and hooves and possibly even the head and tail intact, as illustrated in some of the paintings of the yogi's mat in the \emph{Mysore Yoga Manuscript}. The yogi also wears a patchwork robe made of rags and blankets, carries a rhinoceros' horn as a water pot, which seems rather exotic, and wears a twisted mass of matted locks free of adornments. Figure~\ref{introfig5} presents an AI-generated reconstruction based closely on \emph{Koraṇṭakapaddhati~III}'s description of the yogi and rendered in the style of an eighteenth-century South Indian gouache painting. The\par}
\end{eng}
\end{multicols*}


\begin{center}
    % 1. Measure image width at full height
    \sbox0{\includegraphics[height=\textheight, keepaspectratio]{intro/05.png}}%
    \def\captionwidth{0.2\textwidth}%
    \def\gap{15pt}%
    
    % 2. CAPTION ON THE LEFT (in a zero-width container)
    \makebox[0pt][r]{% Pushes content to the left of the image edge
        \begin{minipage}[b][\textheight][b]{\captionwidth}%
            \captionsetup{justification=raggedleft, singlelinecheck=false, skip=0pt}
            \captionof{figure}{\selectlanguage{english}\setstretch{1.4} Koraṇṭaka's Yogi. This is an AI generated image based on prompts using the unique textual description of a yogi given in \emph{Koraṇṭakapaddhati~III}. (ChatGPT; August 2026).}%
            \vspace{0pt}% Lock bottom alignment
        \end{minipage}%
        \hspace{\gap}%
    }%
    % 3. CENTERED FIGURE
    \begin{minipage}[b][\textheight][b]{\wd0}
        \usebox0%
        \vspace{0pt}% Lock bottom alignment
    \end{minipage}%
\end{center}


\begin{multicols*}{2}
\begin{eng}
\noindent
image is intended as an interpretive reconstruction rather than a historical representation.

The stringed instrument (\emph{tantrīvādya}) and the book (\emph{pustaka}) listed among the yogi's accoutrements are less commonly associated with ascetics and may instead be emblematic of yogis within Koraṇṭaka's tradition. The term \emph{tantrīvādya} appears to correspond to \emph{tatavādya}, a category that includes a range of stringed instruments.\endnote{See the entry on \emph{vādya} in the \emph{Vācaspatyam.}} Although ascetics are known to play such instruments in ritual performance and devotional singing, the use of musical instruments is not attested in earlier works on Haṭhayoga.\endnote{For a reference to the use of the \emph{tantrīvādya} in early Śaiva ritual, see Chemburkar 2022: 197. On the use of the \emph{kiṅgrī} by the Nātha Siddhas, see Mallinson 2021: 87.} However, \emph{Koraṇṭakapaddhati}~\emph{I} prescribes singing the names of the deity for purifying speech and for overcoming the adverse effects of inserting stalks into the urethra during the practice of \emph{vajrolīmudrā}. It is possible that this devotional singing involved the use of an instrument. The description of the yogi's consorts in \emph{Koraṇṭakapaddhati~II} as singers capable of playing a stringed instrument further suggests that devotional singing played an important role in this tradition of Haṭhayoga.

Likewise, books are rarely mentioned among the accoutrements of a yogi, yet they are included in \emph{Koraṇṭakapaddhati~III}'s description.\endnote{This is based on the emendation of \emph{mastakaṃ} to \emph{pustakaṃ} in \emph{Koraṇṭakapaddhati}~\emph{III}. The word \emph{mastakaṃ} is a likely dittographic error as it precedes \emph{mastake} in verse 15.} Another notable exception occurs in a list found in the \emph{Navanāthanavaka}.\endnote{For the reference, see Mallinson 2021: 66, n. 3.} However, several primary sources suggest that books were indeed used by yogis. In one symbolic formulation, the goddess is said to be present as the yogi's guru wherever the book is kept.\endnote{\emph{Khecarīvidyā} 1.27--28. Also see \emph{Haṭhapradīpikā}, 6-chapter version (RORI, ms.\ no.~6756340),~5.107--108, and \emph{Haṃsavilāsa} 4.54.} Although Koraṇṭaka's manuals do not explicitly prescribe the use of books, \emph{Koraṇṭakapaddhati}~\emph{I} refers to the scriptures of Matsyendra, Gorakṣa, Kāki, Kāpālika, and other \emph{siddha}s, as well as to Vaiṣṇava Tantras. This raises the possibility that such works circulated in written form within Koraṇṭaka's tradition.

Despite its emphasis on \emph{vajrolīmudrā} and its portrayal of the yogi in ascetic garb, \emph{Koraṇṭakapaddhati~III} does not reject brahmanical orthodoxy. On the contrary, it praises the use of vedic \emph{mantra}s in worship and enjoins the feeding of Brahmins. It concludes with a unique ritual to be performed by a yogi perfected by the practice of \emph{vajrolīmudrā.} The ritual requires two gold chains, two rings, and a sacrificial gold brick with twelve facets, each bearing a syllable of the Vāsudeva \emph{mantra}. Having smeared the brick with fragrant substances, the yogi worships it and recites the twelve-syllable Vāsudeva \emph{mantra}. Although the use of a gold brick may hark back to the gold bricks of sacrificial altars described in Epic literature,\endnote{For example, in the \emph{Mahābhārata} (14.90.30), Bhīma had an altar built of golden bricks.} the expensive implements required for this rite and the promised worldly outcomes, which include destroying enemies, increasing friends and wealth, and subjugating kings, suggest that these yogis may have enjoyed the patronage of royalty or wealthy nobles who would have benefited from ritual practices that furthered their political ambitions.
\vskip 18pt

\Subsection{The Historical Significance}\label{introsec2.5}
\vskip 9pt

\noindent
The three \emph{Koraṇṭakapaddhati}s of the \emph{Mysore Yoga Manuscript} provide the only surviving account of a tradition of Haṭhayoga attributed to the \emph{siddha} Koraṇṭaka. Likely originating within a Maharashtrian lineage of \emph{siddha}s, these \emph{paddhati}s had reached Mysore by the early nineteenth century. The practice of \emph{āsana} documented in the first \emph{paddhati} was subsequently revised and incorporated into the royal compendium, the \emph{Śrītattvanidhi,} in a section that reflects the prominence of \emph{yogāsana}s in the physical culture of nineteenth-century Mysore. Koraṇṭaka's \emph{paddhati}s therefore represent a tradition of Haṭhayoga that was transmitted from a Maharashtrian ascetic milieu to the royal court of Mysore. The transmission of the first \emph{paddhati} across different social contexts is further evinced by its survival in both a crudely scribed notebook held by a library in Maharashtra (see Figure~\ref{introfig6}) and a beautifully illustrated manuscript produced for the royal court of Mysore.


{\parfillskip=0pt
Read together, the \emph{paddhati}s of Koraṇṭaka constitute an unprecedented account of a tradition of Haṭhayoga that differs from other known traditions in several important respects. Firstly, central to Koraṇṭaka's Haṭhayoga was the practice of \emph{vajrolīmudrā}. The practice had become a highly elaborate discipline requiring extensive preliminary training, including mastery of a strenuous sequence of \emph{āsana}s, intensive cleansing techniques, and methods\par} 
\newpage
\vspace*{\fill}


\noindent
\begin{minipage}{\columnwidth}%
    % 1. Store image inside the minipage scope
    \savebox0{\includegraphics[width=0.65\columnwidth, keepaspectratio]{intro/06.jpg}}%
    \hfill % Right align the image
    % 2. CAPTION ON THE LEFT (anchored to bottom of image height \ht0 + \dp0)
    \makebox[0pt][r]{%
        \begin{minipage}[b][\dimexpr\ht0+\dp0\relax][b]{0.3\columnwidth}
            \captionsetup{justification=raggedleft, singlelinecheck=false, skip=0pt}%
            \captionof{figure}{\selectlanguage{english}\setstretch{1.4} Ms.~No.~29,~2171, f.~2\emph{r}, at the Bhārata Itihāsa Saṃśodhaka Maṇḍala, Pune, Maharashtra.}%
            \par\vspace{0pt}%
        \end{minipage}%
        \hspace{15pt}% Gap between caption and image
    }%
    % 3. PRINT THE IMAGE
    \usebox0%
\end{minipage}
\vspace*{\fill}
\columnbreak

\noindent
for counteracting adverse side effects. It also required considerable resolve on the part of the yogi to endure intense pain, as well as the assistance of others, including sixteen women willing to engage in amorous play and sexual intercourse, and access to a range of materials and resources, from medicinal plants and precious metals to several different types of huts. The practice aimed not only to preserve and raise semen upwards but also to cultivate extreme indifference (\emph{vairāgya}) towards women as objects of sexual desire.

%~ \begin{center}
%~ \includegraphics[width=0.4\columnwidth]{intro/06.jpg}\\
%~ {Figure 6: Outside cover of Ms. No. 29, 2171 at the Bhārata Itihāsa Saṃśodhaka Maṇḍala, Pune.}
%~ \refstepcounter{figure}\label{introfig6}
%~ \end{center}


In a half-verse resembling a Middle Indic proverb, \emph{Koraṇṭakapaddhati}~\emph{II} (3cd) states that \emph{vajrolī} belongs to Viṣṇu, whereas mundane sex belongs to Śiva. Beyond affirming the Vaiṣṇava orientation of Koraṇṭaka's tradition and underscoring the importance of \emph{vajrolī} within it, this statement also points to the earlier history of this \emph{mudrā} in Vaiṣṇava ascetic orders. One of the earliest and most detailed surviving accounts of \emph{vajrolī} appears in the \emph{Dattātreyayogaśāstra}, a thirteenth-century Vaiṣṇava work on Haṭhayoga that also teaches the variants \emph{sahajolī} and \emph{amarolī} found in Koraṇṭaka's tradition. By contrast, early Śaiva works on Haṭhayoga omit \emph{vajrolī}; evidence for the practice of this \emph{mudrā} within Śaivism does not appear until the fifteenth century, in texts such as the \emph{Haṭhapradīpikā} and \emph{Śivasaṃhitā}.

Koraṇṭaka teaches several unprecedented ancillary techniques associated with the practice of \emph{vajrolī}, perhaps the most unusual of which are called the Pigeon Throat, Frog Belly, and Cup. Taught in \emph{Koraṇṭakapaddhati}~\emph{I} and~\emph{II}, these practices involve blowing air through a urethral tube, contracting and relaxing the perineum, drawing a gem in through a urethral tube and expelling it, and drawing air from the tube up through the navel. Among their various benefits, these practices are said to bestow the strength of a lizard. The lizard's strength appears to refer to the yogi's ability to grip the muscles of the pelvic floor strongly. The legendary grip strength of the Bengal monitor lizard was well known in Maharashtra, where it features in the popular story of Chatrapati Shivaji Maharaj. According to the folktale, Shivaji's general Tanaji Malusare used a monitor lizard with ropes attached to it to scale the walls of Kondana Fort during the Battle of Sinhagad (Krishna 2018: 166--167).\endnote{The episode with the monitor lizard features in Tulsidās' ballad (\emph{povāḍā}) on Tanaji Malusare (\emph{cauka} 33). See Acworth 1894: 34--35. We wish to thank Dominic Goodall for bringing this folktale to our attention.}

Koraṇṭaka's yoga tradition also developed the most sophisticated \emph{āsana} practice known to have been documented before the twentieth century. Its one hundred and fourteen postures are divided into six groups and were likely practised in sequences, two of which employ supine and prone positions as linking or transitional positions. Many of the postures involve repetitive movement, while others require extreme spinal flexion and considerable strength. The collection also includes an unprecedented group of rope postures, a pose performed against a wall, and another requiring the assistance of a second person, as depicted by the artist's accompanying illustration (see Section~\ref{introsec3.3}).

The importance of the \emph{āsana} practice in Koraṇṭaka's yoga tradition is articulated differently in each \emph{paddhati}. According to \emph{Koraṇṭakapaddhati}~\emph{I}, the \emph{āsana} practice prepares the yogi for the six cleansing techniques (\emph{ṣaṭkarma}) and makes the body firm and strong (\emph{śārīradārḍhya}). \emph{Koraṇṭakapaddhati}~\emph{II} states that the physical difficulty of \emph{āsana} practice purifies the body's channels (\emph{nāḍī}), whereas \emph{Koraṇṭakapaddhati}~\emph{III} teaches that \emph{āsana} alleviates pain, remedies any imbalances caused by yoga, and enables the yogi to cure his diseases without the help of a doctor. More broadly, in \emph{Koraṇṭakapaddhati}~\emph{II}, Koraṇṭaka's strenuous and dynamic \emph{āsana}s function as preparation for an extreme form of sexual yoga that would have required sufficient fitness and endurance to sustain daily sexual intercourse with sixteen women. \emph{Āsana} appears to play the same role in \emph{Koraṇṭakapaddhati}~\emph{III}. Thus, whereas the relationship between the \emph{āsana} practice and \emph{vajrolī} remains indirect in \emph{Koraṇṭakapaddhati~I}, since one can only infer that \emph{āsana}s support \emph{mudrā}s through the auxiliary system, \emph{Koraṇṭakapaddhati}~\emph{II} and \emph{III} indicate the direct preparatory role of \emph{āsana}s in the practice of \emph{vajrolī}.

It is also of significance that Koraṇṭaka's \emph{paddhati}s are the only surviving account of how the behavioural codes (\emph{yama}) and observances (\emph{niyama}) were understood specifically by proponents of Haṭhayoga. Although commentaries on \emph{yama} and \emph{niyama} survive in other yoga traditions, particularly Pātañjalayoga, little is known about how these guidelines were adapted for Haṭhayoga practitioners, perhaps because most Haṭha texts do not discuss them. Particularly noteworthy is that, in Koraṇṭaka's tradition, celibacy consists in retaining semen while drawing up female sexual fluids rather than in sexual abstinence; hospitality is understood as kindness toward those who have practised the methods of Haṭhayoga rather than as the honouring of guests; and asceticism is defined as carrying out one's religious duty rather than as engaging in penance or bodily mortification.

Koraṇṭaka's tradition is further distinguished by the way it reinterprets techniques known from earlier traditions of Haṭhayoga. A notable example is its interpretation of the three locks (\emph{bandha}): \emph{mūla}, \emph{uḍḍīyāna} and \emph{jālandhara.} These \emph{bandha}s are described as methods of moving the breath through the lower, middle, and upper regions of the torso, respectively, before performing breath retentions (\emph{prāṇāyāma}). In earlier traditions, however, these same locks are explained in terms of contracting or expanding specific parts of the body, such as contracting the anus (\emph{mulabandha}) or stretching the navel region in towards the spine (\emph{uḍḍīyānabandha}).

In other instances, Koraṇṭaka presents variations on standard Haṭhayoga techniques that are unattested in earlier works. \emph{Koraṇṭakapaddhati}~\emph{II}, for example, describes a version of the breathing technique known as the Bellows (\emph{bhastrā}), in which the yogi assumes the Boar Pose (\emph{varāhāsana}) and repeatedly draws air in and out of the anus, like a blacksmith's bellows, hundreds and even thousands of times. Later in the same \emph{paddhati}, holding the breath (\emph{vāyor grahaṇam}) is described as blowing into the urethral tube of the penis and retaining the breath within it.

It is exceptionally rare to possess three manuals documenting a distinct tradition of Haṭhayoga, each presenting a markedly different perspective. Taken together, they illuminate a living tradition of yoga whose unusual and often extreme practices would have set its practitioners apart from other contemporary exponents of Haṭhayoga. Although literature on Haṭhayoga flourished in the early modern period, most surviving works are scholarly compilations that synthesise physical yoga with the more philosophical traditions of Vedānta and Pātañjalayoga. In contrast to this, Koraṇṭaka's \emph{paddhati}s are manuals intended for practitioners and appear to document practices that were in use. Their detailed exposition of \emph{āsana} practice was instrumental in the transmission of Koraṇṭaka's repertoire of postures to the Royal House of Mysore, where aspects of its dynamic style of practice were preserved and adapted before contributing to the dynamic styles of postural yoga that are now practised throughout the world.
\end{eng}
\end{multicols*}

\newpage
~\vfill
\begin{eng}
\centerline{\Section{The Art of \emph{Koraṇṭakapaddhati I}}\label{introsec3}}
\end{eng}
\vskip 22pt


\begin{eng}
\centerline{\Section{From Mysore to the World}\label{introsec4}}
\end{eng}
\vskip 22pt

\centerline{{\small To be added...}}

~\vfill


\newpage

% Figures 7--21 are listed in the Word document but belong to the missing Sections 3 and 4; exact in-text placement cannot be recovered from this source.
\null\vskip 0.1cm
\begin{eng}
\setcounter{section}{4}
\centerline{\Section{Conventions of the Edition and Translation}\label{introsec5}}
\end{eng}
\vskip 22pt

\begin{multicols*}{2}
\begin{eng}
\lettrine[lines=3,lhang=0.04,nindent=0em]{T}{his volume presents a scholarly facsimile edition} of Koraṇṭaka's three \emph{paddhati}s preserved in the Royal House of \emph{Mysore Yoga Manuscript}, comprising high-resolution colour facsimiles, a critically edited transcription of the Sanskrit text in Devanagari script, an annotated English translation, and a Kannada translation.

The illustrated folios are presented differently from those containing only text. Rather than reproducing each illustrated folio in its entirety, enlarged details have been selected to highlight the finer features of the individual \emph{āsana} paintings. Each enlarged illustration is accompanied by the corresponding Devanagari text, followed by the English and Kannada translations. Photographs of the complete illustrated folios are reproduced in Appendix 1.

The Sanskrit edition in Devanagari script has been prepared by correcting and emending, where possible, scribal errors and other textual defects. An endnote in the Devanagari text marks each corrected or emended word, compound, or sentence. The endnotes report the original readings, together with discussions of the emendations and other textual issues, wherever such information is deemed helpful. In certain cases, readings from another manuscript of \emph{Koraṇṭakapaddhati~I}, held at the Bhārata Itihāsa Saṃśodhaka Maṇḍala (ms. no. 29, 2171) in Pune, have been consulted and reported in the endnotes where relevant.\endnote{Neither the Mysore nor the Pune manuscript can be regarded as the exemplar of the other. Overall, the Mysore manuscript contains considerably fewer scribal errors than the Pune manuscript. Nevertheless, each manuscript has its own textual faults and defects, for which the other can, in many instances, supply corrective evidence. For an edition and annotated translation, along with an introduction, of the text of the Pune manuscript (published under the title \emph{Haṭhābhyāsapaddhati}), see Jason Birch, Mark Singleton, and James Mallinson 2025.} Throughout the notes to this edition, the \emph{Mysore Yoga Manuscript} is designated M, and the Pune manuscript P. The Pune manuscript, however, has not been reported in full.

In the Sanskrit edition, certain unconventional spellings and forms of \emph{sandhi} have been tacitly standardised. For example, systematic unconventional spellings such as \emph{haṭa}, \emph{tatva,} and \emph{vindu}, have been standardised as \emph{haṭha}, \emph{tattva,} and \emph{bindu} respectively. Clear spelling errors, such as \emph{asana} for \emph{āsana}, have been corrected, with the original readings recorded in an endnote. We have also standardised certain unconventional instances of \emph{sandhi}, including the use of an \emph{anusvāra} before vowels (e.g., °\emph{prāpakaṃ} \emph{anudvega}°) and the occurrence of unvoiced consonants before voiced consonants (e.g., \emph{kuryāt bindum}). However, \emph{visarga}s appearing before vowels, semivowels, and voiced consonants have been retained (e.g., \emph{bhāvaḥ asaṃcayaḥ}). We have likewise retained unconventional \emph{sandhi} within compounds, such as unresolved vowels between constituent words (e.g., \emph{rajaādyākarṣaṇaṃ}).

The scribe's punctuation has been modified to reflect our understanding and translation of the text. Where necessary, \emph{visarga}s and \emph{daṇḍa}s have been adjusted accordingly. A single \emph{daṇḍa} is used as a full stop, while a double \emph{daṇḍa} marks the end of salutations and verses, follows headings, and is placed before and after the numbers of verses and \emph{āsana} descriptions.

The metrical peculiarities of \emph{Koraṇṭakapaddhati~II} and \emph{III} warrant some comment. Although the Sanskrit register is relatively simple, the authors embellished both texts with irregular but permitted metrical patterns, known as \emph{vipulā}s, most of which are correctly employed, suggesting a reasonable degree of metrical competence.\endnote{\emph{Koraṇṭakapaddhati~II}: \emph{ra-vipulā} occurs in 8c, 11c, 17c, 20c, 33a, 34c, 39a, 55a; \emph{ma-vipulā} in 12a, 18c, 21a, 21c, 48a; \emph{na-vipulā} in 13c, 19a, 24a, 32c, 38a, 43a. \emph{Koraṇṭakapaddhati~III}: \emph{ra-vipulā} occurs in 6e, 7c, 14c, 19c, 20c; \emph{ma-vipulā} in 6c, 9c, 10c, 12c; \emph{bha-vipulā} in 9a, 24c.} There are also instances in which grammatical correctness has been sacrificed to preserve the metre, including the likely deliberate omission of case endings.\endnote{Examples include \emph{Koraṇṭakapaddhati}~\emph{II}, 3cd (case endings have been omitted to preserve metre); 21b (correcting \emph{keli} to \emph{keliḥ} would be unmetrical); 33 (the incorrect form \emph{yāmitavyāḥ,} instead of yamitavyāḥ, has been adopted because it is metrical here and the correct form would not be) and 54 (\emph{vāyu} is left without the \emph{anusvāra} to avoid the unmetrical long on the 5th syllable).} The author's metrical competence also extended to using some sophisticated metres, including \emph{āryā}, although the \emph{āryā} verses have not been well preserved.\endnote{In \emph{Koraṇṭakapaddhati}~\emph{II}, 7 = \emph{upajāti}; 27 = \emph{gīti} (restored with a conjecture); 42 = \emph{gīti}, 54 = \emph{upagīti}. The metre of the following verses has been lost: 37, first half = \emph{gīti}, second half = unmetrical; 50, first three \emph{pāda}s = \emph{sālinī}, fourth \emph{pāda} unmetrical.} Nonetheless, there are also metrical ambiguities and faults that have not been repaired.\endnote{In \emph{Koraṇṭakapaddhati}~\emph{II}, examples of metrical ambiguities include 3a, where strictly speaking, the first verse quarter is not metrical. However, the author may have intended the fifth syllable (\emph{na-s}) to be taken as short; 5c. is unmetrical but sometimes \emph{vajrolī} can be spelt \emph{vajroli}, which may have been intended here without a \emph{visarga;} and 9cd, whose 6th syllable should be long. This could be fixed by emending all three instances of \emph{sahasra} to \emph{sāhasra} (as found in 10c). Metrical Faults in \emph{Koraṇṭakapaddhati}~\emph{II} include 8b (unmetrical as the 2nd, 3rd, 4th syllables = glg); 8ef (hypermetrical); 10a (5th is heavy with no \emph{vipulā} pattern); 19cd (5th is heavy with no \emph{vipulā} pattern); 21ef (ends with gglg); 22ab (5th is heavy with no \emph{vipulā} pattern); 28cd (first half is \emph{anuṣṭubh}, but second half of the verse has lost its metre) and 55e (unmetrical as all syllables are heavy). In \emph{Koraṇṭakapaddhati}~\emph{III}, 1a (5th is heavy with no \emph{vipulā} pattern); 5c (5th is long with no \emph{vipulā} pattern); 16b (syllables 2-4 are glg); 16c (syllables 2-3 are both light) and 24a (5th is heavy with no \emph{vipulā} pattern).} Both \emph{paddhati}s also contain prose insertions that do not appear to be original and may derive from marginal notes that were subsequently incorporated into the main text.

The English translation errs on the side of a literal rendering of the Sanskrit edition. Words supplied to make the English fluent and grammatically correct are enclosed in square brackets. Sanskrit words that have not been translated but are likely to be familiar to readers, such as \emph{āsana} and \emph{prāṇa}, are italicised and transliterated according to IAST conventions. The Kannada translation closely follows the English. In some cases, Sanskrit terms familiar to Kannada readers have been retained, particularly the names of \emph{āsana}s and other techniques.
\end{eng}
\end{multicols*}

\newpage

\begin{multicols*}{2}
\begin{eng}
\theendnotes
\end{eng}
\end{multicols*}
\end{nohyp}
